\documentclass{beamer}

\usetheme{Madrid}

%================= LANGUAGE =================
\usepackage[utf8]{inputenc}
\usepackage[T5]{fontenc}
\usepackage[vietnamese]{babel}

%================= PACKAGE =================
\usepackage{amsmath}
\usepackage{tikz}
\usepackage{xcolor}
\usepackage{tcolorbox}
\usepackage{fontawesome5}

%================= COLOR =================
\definecolor{main1}{RGB}{102,126,234}
\definecolor{main2}{RGB}{118,75,162}

\setbeamercolor{frametitle}{fg=white,bg=main1}
\setbeamercolor{title}{fg=white,bg=main2}

%================= BACKGROUND =================
\setbeamertemplate{background canvas}{
\begin{tikzpicture}[remember picture,overlay]
\path[left color=main1!20,right color=main2!20]
(current page.south west) rectangle (current page.north east);
\end{tikzpicture}
}

%================= BOX =================
\tcbset{
  colback=white,
  colframe=main1,
  arc=4mm,
  boxrule=1pt,
  drop shadow
}

%================= INFO =================
\title{Biểu diễn đồ thị}
\author{Nhóm 2}
\date{\today}

\begin{document}

\frame{\titlepage}

%--------------------------------
\begin{frame}{Định nghĩa}

\begin{tcolorbox}
\textbf{\faProjectDiagram Biểu diễn đồ thị}

Biểu diễn đồ thị là cách lưu trữ đồ thị $G = (V,E)$.

\begin{itemize}
\item[\faCircle] $V$: tập đỉnh
\item[\faLink] $E$: tập cạnh
\end{itemize}
\end{tcolorbox}

\vspace{0.2cm}

\begin{tcolorbox}
\textbf{\faBullseye Mục tiêu}
\begin{itemize}
\item[\faCheckCircle] Tiết kiệm bộ nhớ
\item[\faBolt] Duyệt nhanh
\item[\faCogs] Phù hợp thuật toán
\end{itemize}
\end{tcolorbox}

\end{frame}

%--------------------------------
\begin{frame}{Đồ thị ví dụ chung}

\begin{columns}

\column{0.5\textwidth}

\begin{tikzpicture}[scale=0.9]
\node[circle,draw,fill=main1!20,thick] (1) at (0,0) {1};
\node[circle,draw,fill=main2!20,thick] (2) at (2,1) {2};
\node[circle,draw,fill=main2!20,thick] (3) at (2,-1) {3};
\node[circle,draw,fill=main1!20,thick] (4) at (4,0) {4};

\draw[thick] (1)--(2);
\draw[thick] (1)--(3);
\draw[thick] (2)--(4);
\draw[thick] (3)--(4);
\end{tikzpicture}

\column{0.5\textwidth}

\begin{tcolorbox}
\[
V=\{1,2,3,4\}
\]
\[
E=\{(1,2),(1,3),(2,4),(3,4)\}
\]
\end{tcolorbox}

\end{columns}

\end{frame}

%--------------------------------
\begin{frame}{Danh sách cạnh (Edge List)}

\begin{columns}

\column{0.5\textwidth}

\begin{tikzpicture}[scale=0.9]
\node[circle,draw,fill=main1!20,thick] (1) at (0,0) {1};
\node[circle,draw,fill=main2!20,thick] (2) at (2,1) {2};
\node[circle,draw,fill=main2!20,thick] (3) at (2,-1) {3};
\node[circle,draw,fill=main1!20,thick] (4) at (4,0) {4};

\draw[thick] (1)--(2);
\draw[thick] (1)--(3);
\draw[thick] (2)--(4);
\draw[thick] (3)--(4);
\end{tikzpicture}

\column{0.5\textwidth}

\begin{tcolorbox}
\textbf{\faProjectDiagram Biểu diễn}
\[
(1,2), (1,3), (2,4)
\]
\end{tcolorbox}

\vspace{0.3cm}

\begin{tcolorbox}
\textbf{\faInfoCircle Đặc điểm}
\begin{itemize}
\item Dễ lưu, dễ nhập
\item Không tối ưu khi duyệt
\end{itemize}
\end{tcolorbox}

\end{columns}

\end{frame}

%--------------------------------
\begin{frame}{Ma trận kề}

\begin{columns}

\column{0.5\textwidth}

\begin{tikzpicture}[scale=0.9]
\node[circle,draw,fill=main1!20,thick] (1) at (0,0) {1};
\node[circle,draw,fill=main2!20,thick] (2) at (2,1) {2};
\node[circle,draw,fill=main2!20,thick] (3) at (2,-1) {3};
\node[circle,draw,fill=main1!20,thick] (4) at (4,0) {4};

\draw[thick] (1)--(2);
\draw[thick] (1)--(3);
\draw[thick] (2)--(4);
\draw[thick] (3)--(4);
\end{tikzpicture}

\column{0.5\textwidth}

\begin{tcolorbox}
\[
\begin{bmatrix}
0 & 1 & 1 & 0 \\
1 & 0 & 0 & 1 \\
1 & 0 & 0 & 1 \\
0 & 1 & 1 & 0
\end{bmatrix}
\]
\end{tcolorbox}

\vspace{0.3cm}

\begin{tcolorbox}
\textbf{\faInfoCircle Đặc điểm}
\begin{itemize}
\item Kiểm tra cạnh nhanh O(1)
\item Tốn bộ nhớ O($n^2$)
\end{itemize}
\end{tcolorbox}

\end{columns}

\end{frame}

%--------------------------------
\begin{frame}{Danh sách kề}

\begin{columns}

\column{0.5\textwidth}

\begin{tikzpicture}[scale=0.9]
\node[circle,draw,fill=main1!20,thick] (1) at (0,0) {1};
\node[circle,draw,fill=main2!20,thick] (2) at (2,1) {2};
\node[circle,draw,fill=main2!20,thick] (3) at (2,-1) {3};
\node[circle,draw,fill=main1!20,thick] (4) at (4,0) {4};

\draw[thick] (1)--(2);
\draw[thick] (1)--(3);
\draw[thick] (2)--(4);
\draw[thick] (3)--(4);
\end{tikzpicture}

\column{0.5\textwidth}

\begin{tcolorbox}
\begin{itemize}
\item 1: 2, 3
\item 2: 1, 4
\item 3: 1, 4
\item 4: 2, 3
\end{itemize}
\end{tcolorbox}

\vspace{0.3cm}

\begin{tcolorbox}
\textbf{\faInfoCircle Đặc điểm}
\begin{itemize}
\item Tiết kiệm bộ nhớ
\item Duyệt nhanh
\end{itemize}
\end{tcolorbox}

\end{columns}

\end{frame}

%--------------------------------
\begin{frame}{So sánh 3 cách}

\begin{columns}

\column{0.5\textwidth}
\begin{tcolorbox}
\textbf{\faBalanceScale Edge List}
\begin{itemize}
\item Lưu đơn giản
\item Không phù hợp BFS/DFS
\end{itemize}
\end{tcolorbox}

\vspace{0.2cm}

\begin{tcolorbox}
\textbf{\faBalanceScale Adjacency Matrix}
\begin{itemize}
\item Kiểm tra cạnh nhanh nhất
\item Chậm khi duyệt toàn bộ
\end{itemize}
\end{tcolorbox}

\column{0.5\textwidth}

\begin{tcolorbox}
\textbf{\faBalanceScale Adjacency List}
\begin{itemize}
\item Duyệt nhanh nhất
\item Dùng phổ biến nhất
\end{itemize}
\end{tcolorbox}

\end{columns}

\end{frame}

%--------------------------------
\begin{frame}{Khi nào nên dùng?}

\begin{columns}

\column{0.5\textwidth}

\begin{tcolorbox}
\textbf{\faLightbulb Edge List}
\begin{itemize}
\item Khi chỉ cần lưu cạnh
\item Bài toán đơn giản
\end{itemize}
\end{tcolorbox}

\vspace{0.2cm}

\begin{tcolorbox}
\textbf{\faLightbulb Adjacency Matrix}
\begin{itemize}
\item Đồ thị dày ($m \approx n^2$)
\item Cần kiểm tra nhanh có cạnh hay không
\end{itemize}
\end{tcolorbox}

\column{0.5\textwidth}

\begin{tcolorbox}
\textbf{\faLightbulb Adjacency List}
\begin{itemize}
\item Đồ thị thưa
\item Dùng BFS, DFS
\item Thi lập trình
\end{itemize}
\end{tcolorbox}

\end{columns}

\end{frame}

\end{document}
